\subsection{P4 动态滚动重排结果}

本节基于 \verb|data/processed/p4| 的最新输出，对 P4 动态重排结果进行验证和解释。实验读取 12 个 P3 基准任务和 3 个动态事件，其中 2 个事件直接影响计划。P4 采用“可行性优先、服务质量其次”的评价顺序：轨迹冲突、封锁区违规、AMR 延误违规、缺失路径和电量阈值违反必须全部为 0；在此基础上，再比较延期、扰动任务数、$C_{\max}$、$F_2$ 和 $F_3$。最终滚动时域重排结果如下：

\begin{table}[H]
\centering
\small
\begin{tabular}{p{0.42\linewidth}r}
\toprule
指标 & 数值 \\
\midrule
动态事件数 & 3 \\
受影响事件数 & 2 \\
新增动态任务数 & 1 \\
变更或插入任务数 & 4 \\
P3 基准 $C_{\max}$ & 46.596970 \\
P4 新 $C_{\max}$ & 61.918082 \\
P4 总延期 \verb|total_delay| & 23.405971 \\
相对基准正向完工偏移 \verb|TotalAddedDelay| & 43.996446 \\
剩余轨迹冲突 & 0 \\
封锁区违规 & 0 \\
AMR 延误违规 & 0 \\
缺失路径数 & 0 \\
电量阈值违反数 & 0 \\
F2 & 2182.068048 \\
F3 & 507.179375 \\
\bottomrule
\end{tabular}
\caption{P4 滚动重排主结果}
\label{tab:p4-summary}
\end{table}

表 \ref{tab:p4-summary} 展示动态事件后的最终可行性和服务质量。这里需要区分两个指标：\verb|total_delay| 是按任务时间窗计算的延期总量，数值为 23.405971；\verb|TotalAddedDelay| 是相对 P3/新任务基线的正向完工时间偏移总和，数值为 43.996446。二者含义不同，不能混写。

\subsection{动态事件影响范围}

P4 首先对三个动态事件逐一做影响分析，结果见表 \ref{tab:p4-event-impact}。影响分析把事件分为“触发重排”和“记录但不扰动”两类。E01 定义了封锁区域和封锁时间窗，但最终轨迹 footprint 没有进入该矩形区域，因此影响任务数为 0；E02 是 AMR2 在 $[24,30]$ 内延误，影响 AMR2 的 T05 和 T10；E03 是新增任务 T13，最终插入 AMR4。

\begin{table}[H]
\centering
\scriptsize
\begin{tabular}{p{0.10\linewidth}p{0.16\linewidth}p{0.12\linewidth}p{0.16\linewidth}p{0.28\linewidth}}
\toprule
事件 & 类型 & 目标 & 受影响任务 & 影响解释 \\
\midrule
E01 & \verb|area_block| & BLOCK\_01 & 无 & 封锁窗口内没有轨迹 footprint 与封锁矩形相交 \\
E02 & \verb|amr_delay| & AMR2 & T05, T10 & AMR2 在 $[24,30]$ 延误，T05 与该窗口重叠，T10 属于同车后续尾段 \\
E03 & \verb|new_task| & T13 & T13 & 新任务进入滚动窗口，最终插入 AMR4 \\
\bottomrule
\end{tabular}
\caption{P4 动态事件影响范围}
\label{tab:p4-event-impact}
\end{table}

这一结果表明，P4 采用事件影响驱动的重排触发机制。封锁事件只在轨迹 footprint 与封锁矩形发生时空交叠时触发重排；AMR 延误事件释放与延误窗口相交的任务及同车后续尾段；新增任务事件则进入未来任务池并参与插入位置搜索。这样的处理方式把动态事件转化为明确的重优化对象，使滚动窗口内的调整范围与实际受影响区域保持一致。

\begin{figure}[!htbp]
\centering
\includegraphics[width=0.90\linewidth]{p4_gantt_events.png}
\caption{P4 动态事件与重排甘特图}
\label{fig:p4-gantt}
\end{figure}

\begin{figure}[!htbp]
\centering
\includegraphics[width=0.90\linewidth]{p4_trajectory_map.png}
\caption{P4 重排后轨迹地图}
\label{fig:p4-map}
\end{figure}

\begin{figure}[!htbp]
\centering
\includegraphics[width=0.98\linewidth]{p4_dynamic_process_snapshots.png}
\caption{P4 动态滚动重排过程快照}
\label{fig:p4-process}
\end{figure}

图 \ref{fig:p4-gantt} 从时间轴上展示动态事件与任务重排之间的关系，图 \ref{fig:p4-map} 从仓库平面上展示重排后的轨迹空间分布，图 \ref{fig:p4-process} 进一步把 GIF 动图中的四个关键时刻整理为静态快照。三张图共同展示 P4 对时间表、空间轨迹和动态过程的同步更新。

该实验的关键变化是 AMR2 的 T05 与延误事件 E02 时间窗相交。P4 将 AMR2 的重叠任务及其后续尾段从固定集合中释放，并从延误结束后重新安排，最终得到 \verb|delayed_amr_violation_count=0|。新增任务 T13 被插入 AMR4，并前置于 T11。该安排用 T11 的局部后移换取 AMR2 尾段负担下降，使滚动重排方案在硬约束可行的同时降低延期和完工时间。

\subsection{重排任务明细}

表 \ref{tab:p4-task-change} 给出最终发生变化或新增的任务。P4 共改变或插入 4 个任务：T05、T10、T13 和 T11。其中 T05 和 T10 仍由 AMR2 执行，变化原因来自 E02；T13 是新增任务，被插入 AMR4；T11 原本也是 AMR4 任务，但由于 T13 被前置，T11 被推迟执行。

\begin{table}[H]
\centering
\scriptsize
\resizebox{\linewidth}{!}{%
\begin{tabular}{p{0.08\linewidth}p{0.09\linewidth}p{0.09\linewidth}rrrrrp{0.20\linewidth}}
\toprule
任务 & 原 AMR & 新 AMR & 原开始 & 新开始 & 原完成 & 新完成 & 延后量 & 原因 \\
\midrule
T05 & AMR2 & AMR2 & 21.101333 & 32.939000 & 27.972444 & 39.810111 & 11.837667 & AMR2 延误后重排 \\
T10 & AMR2 & AMR2 & 34.353444 & 51.191111 & 44.840111 & 61.677778 & 16.837667 & AMR2 尾段重排，并加入 5s 冲突等待 \\
T13 & -- & AMR4 & -- & 34.441920 & -- & 43.453031 & -- & 新增任务释放后插入 AMR4 \\
T11 & AMR4 & AMR4 & 35.146465 & 50.467577 & 46.596970 & 61.918082 & 15.321112 & T13 前置后被推迟 \\
\bottomrule
\end{tabular}
}
\caption{P4 重排任务明细}
\label{tab:p4-task-change}
\end{table}

从表 \ref{tab:p4-task-change} 可以看出，P4 同时完成了两类动态调整。E02 改变 AMR2 的可用状态，使 T05 和 T10 从延误窗口之后重新排布；E03 引入新任务 T13，并通过候选插入搜索选择 AMR4。当前解把 T13 放到 AMR4，并让 T11 后移，相当于用一个额外扰动任务换取更低的总延期和更短的系统完工时间。该选择符合 P4 的目标层次：先保持硬约束为 0，再在可行域内降低延期和 $C_{\max}$。

\subsection{可行性验证}

P4 最终方案的硬约束验证结果集中体现在表 \ref{tab:p4-feasibility}。表中的 0 来自最终 \verb|trajectory_schedule.csv| 的轨迹级复核，覆盖 footprint 冲突、封锁区进入、AMR 延误窗口重叠、路径缺失和电量阈值违反。

\begin{table}[H]
\centering
\small
\begin{tabular}{p{0.42\linewidth}r p{0.34\linewidth}}
\toprule
指标 & 数值 & 含义 \\
\midrule
剩余轨迹冲突数 & 0 & 最终轨迹中无 footprint 或节点容量冲突 \\
封锁区违规数 & 0 & 封锁区时间窗内无轨迹进入封锁矩形 \\
AMR 延误违规数 & 0 & AMR2 延误窗口内无任务执行区间重叠 \\
缺失载货路径数 & 0 & 所有载货路径均能在 P1 mixed 路径库中找到 \\
缺失空驶衔接路径数 & 0 & 所有空驶衔接路径均能在 P1 mixed 路径库中找到 \\
电量阈值违反数 & 0 & 所有 AMR 完成任务后电量不低于安全阈值 \\
全部任务已排程 & 1 & 静态任务和动态新增任务均已排入最终计划 \\
\bottomrule
\end{tabular}
\caption{P4 最终方案硬约束验证}
\label{tab:p4-feasibility}
\end{table}

表 \ref{tab:p4-feasibility} 给出最终动态可行性证据。AMR 延误违规数为 0，说明 E02 已被轨迹层完全消化：T05 原计划从 21.101333 执行到 27.972444，与 AMR2 的 $[24,30]$ 延误窗口重叠；P4 将其新开始时间推迟到 32.939000，从而避开延误窗口。T10 作为 AMR2 后续任务，也被相应推迟到 51.191111 开始。最终 AMR2 在不可用窗口内没有任务执行区间。

\begin{figure}[!htbp]
\centering
\includegraphics[width=0.88\linewidth]{p4_event_impact_metrics.png}
\caption{P4 事件影响与可行性指标}
\label{fig:p4-impact}
\end{figure}

\begin{figure}[!htbp]
\centering
\includegraphics[width=0.88\linewidth]{p4_method_comparison.png}
\caption{P4 方法对比输出}
\label{fig:p4-method}
\end{figure}

图 \ref{fig:p4-impact} 将动态事件影响、硬约束违反数和目标函数指标集中展示，图 \ref{fig:p4-method} 则给出重排策略的对照输出。二者共同强调 P4 的评价顺序：先看轨迹冲突、封锁违规和 AMR 延误违规是否清零，再比较延期、扰动任务数和完工时间。这一顺序与实际仓储系统一致，因为不可执行的低成本方案没有运营意义。

\subsection{目标函数指标解释}

P4 的主要目标指标见表 \ref{tab:p4-objective}。\verb|priority_late_count| 和 \verb|late_count| 描述延期任务数量层面的服务质量；\verb|priority_delay| 和 \verb|total_delay| 描述延期幅度；$C_{\max}$ 表示系统最终完工时间；$F_2$ 与 $F_3$ 分别综合时间效率和运行成本。

\begin{table}[H]
\centering
\small
\begin{tabular}{p{0.34\linewidth}r p{0.38\linewidth}}
\toprule
指标 & 数值 & 解释 \\
\midrule
\verb|priority_late_count| & 6.0 & 延期任务按优先级加权后的数量 \\
\verb|late_count| & 3 & 出现正延期的任务数量 \\
\verb|priority_delay| & 46.811941 & 延期量按任务优先级加权后的总和 \\
\verb|total_delay| & 23.405971 & 按任务软时间窗计算的总延期 \\
$C_{\max}$ & 61.918082 & P4 最终系统完工时间 \\
$F_2$ & 2182.068048 & $30\sum_i w_iT_i+20\sum_iT_i+5C_{\max}$ \\
$F_3$ & 507.179375 & 空驶成本、负载不均衡和总能耗综合指标 \\
\bottomrule
\end{tabular}
\caption{P4 目标函数指标}
\label{tab:p4-objective}
\end{table}

动态事件改变了 P3 基准计划的可用时间和任务集合，特别是 AMR2 的不可用窗口压缩了其后续任务的可用区间。P4 先恢复硬可行性，再在可行域内优化服务指标；当前方案通过把 T13 插入 AMR4 来降低 AMR2 尾段负担。最终方案包含 3 个延期任务，但轨迹冲突、封锁违规和 AMR 延误违规全部为 0。

\begin{figure}[!htbp]
\centering
\includegraphics[width=0.96\linewidth]{p4_weight_sensitivity.png}
\caption{P4 目标权重配置敏感性结果}
\label{fig:p4-weight}
\end{figure}

图 \ref{fig:p4-weight} 比较服务等级优先、均衡权重和稳定性优先三类目标配置。在当前算例中，三类配置均选择同一条可行重排方案，说明本算例的主导因素来自硬约束恢复和新增任务插入结构；局部权重变化没有改变最终候选选择。该结果支持本文采用均衡权重作为主报告方案。

\subsection{与中间候选方案的比较}

在调试和优化 P4 时，曾出现两个具有代表性的 AMR2 插入候选。第一个候选把 T13 插入 AMR2 并排在 T05 之前，虽然延期任务数较少，但会让 T05 和 T10 大幅后移，总延期为 41.794333，$C_{\max}=73.331$，$F_2=3710.201667$。第二个候选把 T13 插在 AMR2 的 T05 与 T10 之间，总延期降为 27.105444，但 $C_{\max}$ 仍为 73.241，AMR2 尾段仍然较重。最终方案把 T13 插入 AMR4 并前置于 T11，总延期进一步降到 23.405971，$C_{\max}$ 降到 61.918082。

\begin{table}[H]
\centering
\small
\begin{tabular}{p{0.32\linewidth}rrrr}
\toprule
方案 & 扰动任务数 & 总延期 & $C_{\max}$ & $F_2$ \\
\midrule
AMR2 前插 T13 候选 & 3 & 41.794333 & 73.331000 & 3710.201667 \\
AMR2 中插 T13 候选 & 3 & 27.105444 & 73.241000 & 2596.270556 \\
当前 AMR4 前插 T13 方案 & 4 & 23.405971 & 61.918082 & 2182.068048 \\
\bottomrule
\end{tabular}
\caption{P4 新任务插入候选效果比较}
\label{tab:p4-candidate-compare}
\end{table}

表 \ref{tab:p4-candidate-compare} 表明，当前方案以 1 个额外扰动任务换取了更低的总延期、更短的 $C_{\max}$ 和更优的 $F_2$。相对 AMR2 中插候选，当前方案总延期降低
\[
\frac{27.105444-23.405971}{27.105444}\approx 13.65\%,
\]
$C_{\max}$ 降低
\[
\frac{73.241000-61.918082}{73.241000}\approx 15.46\%.
\]
相对最初 AMR2 前插候选，当前方案总延期降低约 $44.00\%$，$F_2$ 降低约 $41.19\%$。这体现了 P4 短名单修复机制的作用：候选方案先经过轨迹检测和有限等待修复，再按字典序目标比较，因此能够保留初始冲突较多但修复后目标更优的插入位置。

\subsection{图形结果解释}

图 \ref{fig:p4-gantt} 的主要作用是说明时间维度上的重排逻辑。AMR2 的 T05 被推迟到延误窗口之后，T10 作为同车后续任务继续后移；AMR4 在完成 T08 后插入 T13，再执行 T11。这个甘特图可以直接看出 P4 的滚动时域特点：事件发生前的已执行任务保持不变，事件后的未来尾段被局部调整。

图 \ref{fig:p4-map} 则说明空间维度上的可行性。新增任务 T13 从 P4 到 SORT2，插入 AMR4 后没有造成封锁区违规，也没有留下轨迹冲突。轨迹图与甘特图配合使用，可以避免只从时间表判断动态重排是否成功。对多 AMR 仓储系统而言，时间上不重叠并不必然代表空间上安全，空间上绕开也不必然代表时间窗可接受；P4 必须同时满足两者。

图 \ref{fig:p4-impact} 将事件影响、硬约束违反和目标指标放在同一张图里，形成最终方案的可行性证据链。图 \ref{fig:p4-method} 汇总滚动时域主结果，图 \ref{fig:p4-weight} 展示目标权重配置下的指标稳定性。三类图分别对应“事件影响”“方法结果”和“目标敏感性”，使 P4 的实验结论从时间、空间、目标三个层面得到支撑。

\subsection{适用范围与扩展实验}

P4 的实现定位为动态仓储运行中的在线滚动重排启发式。它把完整时空 MILP 中的任务指派、路径选择、等待和轨迹互斥约束转化为滚动窗口内的候选插入、短名单修复和等待冲突修复流程。该实现保留了混合整数规划的决策结构和目标层级，同时把求解范围限制在受影响任务池内，适合事件到达后快速恢复可执行计划。

当前解的数学含义是“给定候选集合与有界修复策略下的高质量动态可行解”。它的优势在于指标清晰、可复现、运行时间可接受，并且最终轨迹通过硬约束检查。在此基础上，可以进一步加入小规模精确模型作为离线上界对照：例如只对受影响 AMR 和相邻时间窗建立局部 MILP，或者把候选插入、等待动作离散化后，用 0-1 规划选择动作组合。

离线横向对比可以把“仅等待”“全局重排”“滚动时域重排”作为三类基线分别运行，并统一比较新增延期、扰动任务数、冲突消解成功率和运行时间。这样可以进一步量化滚动时域方法的收益来源：它在保持历史计划稳定的同时，对受影响未来尾段进行重优化，从而兼顾执行稳定性和动态服务质量。

总体而言，P4 已经完成动态仓储场景下的保底要求和进阶要求：能够识别封锁、延误和新增任务事件；能够释放受影响任务和未来尾段；能够将新增任务插入合适 AMR；能够输出结构化重排表、可行性指标和动态可视化；最终方案在轨迹冲突、封锁区违规和 AMR 延误违规上均为 0，并通过局部搜索将总延期控制在 23.405971。

